\begin{ex}
Tính giá trị của biểu thức: $10 - 3 + 2 - 4$.
\shortans{5}
\loigiai{
Thực hiện lần lượt: $10 - 3 = 7$; $7 + 2 = 9$; $9 - 4 = 5$.
\textbf{Đáp số:} $5$.
}
\end{ex}

\begin{ex}
Tìm số tự nhiên lớn nhất có một chữ số trừ đi số bé nhất có một chữ số.
\shortans{9}
\loigiai{
Số lớn nhất có một chữ số là 9, bé nhất là 0: $9 - 0 = 9$.
\textbf{Đáp số:} $9$.
}
\end{ex}

\begin{ex}
Tìm số tự nhiên $x$ thỏa mãn: $x + 4 = 10 - 2$.
\shortans{4}
\loigiai{
Vế phải: $10 - 2 = 8$.\nTa có: $x + 4 = 8 \Rightarrow x = 8 - 4 = 4$.
\textbf{Đáp số:} $4$.
}
\end{ex}

\begin{ex}
Tìm số tự nhiên $y$ thỏa mãn: $9 - y = 3 + 3$.
\shortans{3}
\loigiai{
Vế phải: $3 + 3 = 6$.\nTa có: $9 - y = 6 \Rightarrow y = 9 - 6 = 3$.
\textbf{Đáp số:} $3$.
}
\end{ex}

\begin{ex}
Bạn An có 5 viên kẹo, Bình có nhiều hơn An 2 viên kẹo, Cúc có ít hơn Bình 3 viên kẹo. Hỏi Cúc có bao nhiêu viên kẹo?
\shortans{4}
\loigiai{
Bình có số kẹo là: $5 + 2 = 7$ (viên kẹo).\nCúc có số kẹo là: $7 - 3 = 4$ (viên kẹo).
\textbf{Đáp số:} $4$.
}
\end{ex}

\begin{ex}
Tổng số kẹo của cả ba bạn An, Bình và Cúc trong bài toán trên là bao nhiêu viên kẹo?
\shortans{16}
\loigiai{
Tổng số kẹo là: $5 + 7 + 4 = 16$ (viên kẹo).
\textbf{Đáp số:} $16$.
}
\end{ex}

\begin{ex}
Một cây nến dài 10 cm, bị đốt cháy đi 6 cm. Hỏi cây nến còn lại dài bao nhiêu xăng-ti-mét?
\shortans{4}
\loigiai{
Độ dài còn lại của cây nến là: $10 - 6 = 4$ (cm).
\textbf{Đáp số:} $4$.
}
\end{ex}

\begin{ex}
Trong một chuồng nuôi thỏ có 6 con thỏ trắng và 4 con thỏ nâu. Bác nông dân bán đi 5 con thỏ trắng. Hỏi trong chuồng còn lại bao nhiêu con thỏ cả hai loại?
\shortans{5}
\loigiai{
Số thỏ trắng còn lại là: $6 - 5 = 1$ (con thỏ trắng).\nTổng số thỏ còn lại trong chuồng là: $1 + 4 = 5$ (con thỏ).
\textbf{Đáp số:} $5$.
}
\end{ex}

\begin{ex}
Tìm số đứng liền trước của số bé nhất có hai chữ số.
\shortans{9}
\loigiai{
Số bé nhất có hai chữ số là 10. Số liền trước của 10 là $9$.
\textbf{Đáp số:} $9$.
}
\end{ex}

\begin{ex}
Tìm số đứng liền sau của số bé nhất có một chữ số khác 0.
\shortans{2}
\loigiai{
Số bé nhất khác 0 là 1. Số liền sau của 1 là $2$.
\textbf{Đáp số:} $2$.
}
\end{ex}

\begin{ex}
Có bao nhiêu số tự nhiên có một chữ số lớn hơn 4?
\shortans{5}
\loigiai{
Các số đó là: $5, 6, 7, 8, 9$. Có tất cả $5$ số.
\textbf{Đáp số:} $5$.
}
\end{ex}

\begin{ex}
Hình tam giác có ít hơn khối lập phương bao nhiêu đỉnh?
\shortans{5}
\loigiai{
Hình tam giác có 3 đỉnh, khối lập phương có 8 đỉnh: $8 - 3 = 5$ (đỉnh).
\textbf{Đáp số:} $5$.
}
\end{ex}

\begin{ex}
Tổng số mặt của một khối lập phương và một hình vuông là bao nhiêu? (Khối lập phương có 6 mặt, hình vuông tính là 1 mặt phẳng).
\shortans{7}
\loigiai{
Tổng số mặt là: $6 + 1 = 7$.
\textbf{Đáp số:} $7$.
}
\end{ex}

\begin{ex}
Bạn Hùng có 10 que tính, dùng 4 que xếp hình vuông và 3 que xếp hình tam giác. Hỏi Hùng còn thừa bao nhiêu que tính?
\shortans{3}
\loigiai{
Tổng số que tính đã dùng: $4 + 3 = 7$ (que tính).\nSố que tính còn thừa là: $10 - 7 = 3$ (que tính).
\textbf{Đáp số:} $3$.
}
\end{ex}

\begin{ex}
Hai số tự nhiên liên tiếp có tổng bằng 9. Số lớn hơn trong hai số đó là số mấy?
\shortans{5}
\loigiai{
Hai số liên tiếp có tổng bằng 9 là 4 và 5. Số lớn hơn là số $5$.
\textbf{Đáp số:} $5$.
}
\end{ex}

\begin{ex}
Hai số tự nhiên liên tiếp có tổng bằng 7. Số bé hơn trong hai số đó là số mấy?
\shortans{3}
\loigiai{
Hai số liên tiếp có tổng bằng 7 là 3 và 4. Số bé hơn là số $3$.
\textbf{Đáp số:} $3$.
}
\end{ex}

\begin{ex}
Trong một hàng dọc có 10 bạn học sinh. Bạn Nam đứng ở vị trí thứ 6 từ đầu hàng. Hỏi phía sau bạn Nam có bao nhiêu bạn?
\shortans{4}
\loigiai{
Số bạn học sinh phía sau Nam là: $10 - 6 = 4$ (bạn).
\textbf{Đáp số:} $4$.
}
\end{ex}

\begin{ex}
Một nhóm học sinh có 6 bạn nữ và 3 bạn nam. Thêm 1 bạn nam vào nhóm. Hỏi lúc này số bạn nữ nhiều hơn số bạn nam bao nhiêu bạn?
\shortans{2}
\loigiai{
Số bạn nam lúc này là: $3 + 1 = 4$ (bạn).\nSố bạn nữ nhiều hơn số bạn nam là: $6 - 4 = 2$ (bạn).
\textbf{Đáp số:} $2$.
}
\end{ex}

\begin{ex}
Tìm số tự nhiên $a$ biết rằng: $a + a = 10$.
\shortans{5}
\loigiai{
Hai số giống nhau có tổng bằng 10 là 5 (vì $5 + 5 = 10$). Vậy $a = 5$.
\textbf{Đáp số:} $5$.
}
\end{ex}

\begin{ex}
Tìm số tự nhiên $b$ biết rằng: $b + b = 6$.
\shortans{3}
\loigiai{
Hai số giống nhau có tổng bằng 6 là 3 (vì $3 + 3 = 6$). Vậy $b = 3$.
\textbf{Đáp số:} $3$.
}
\end{ex}

\begin{ex}
Số điểm kiểm tra học kì 1 môn Toán của bạn Lan là số lớn nhất có một chữ số cộng thêm 1. Hỏi Lan được mấy điểm?
\shortans{10}
\loigiai{
Số lớn nhất có một chữ số là 9: $9 + 1 = 10$ (điểm 10).
\textbf{Đáp số:} $10$.
}
\end{ex}

\begin{ex}
Tính tổng các số chẵn có một chữ số khác 0 và bé hơn 6.
\shortans{6}
\loigiai{
Các số chẵn khác 0 bé hơn 6 là 2 và 4.\nTổng là: $2 + 4 = 6$.
\textbf{Đáp số:} $6$.
}
\end{ex}

\begin{ex}
Tính tổng các số lẻ có một chữ số bé hơn 6.
\shortans{9}
\loigiai{
Các số lẻ bé hơn 6 là: 1, 3, 5.\nTổng là: $1 + 3 + 5 = 9$.
\textbf{Đáp số:} $9$.
}
\end{ex}

\begin{ex}
Có 10 quả cam, chia đều cho hai bạn Mai và Lan. Hỏi mỗi bạn nhận được bao nhiêu quả cam?
\shortans{5}
\loigiai{
Chia đều 10 quả cam cho hai bạn: mỗi bạn nhận được $10 : 2 = 5$ (quả cam).
\textbf{Đáp số:} $5$.
}
\end{ex}

\begin{ex}
Cho dãy số: $1, 4, 7, 10$. Số đứng ở vị trí thứ ba trong dãy là số nào?
\shortans{7}
\loigiai{
Dãy số gồm: số thứ nhất là 1, số thứ hai là 4, số thứ ba là 7, số thứ tư là 10.\nSố đứng ở vị trí thứ ba là số $7$.
\textbf{Đáp số:} $7$.
}
\end{ex}
